% ECONOMICS_PROBLEM: {"id": "C72-14", "title": "The Catch-Up optimal-draw conjecture for consecutive integers", "jel_code": "C72", "source_id": "OP-0589"}
% FIRST_POSTED: 2026-10-09
% ATTEMPT_STATUS: unresolved
% ENTRY_KIND: research_attempt
\documentclass[11pt,a4paper]{article}

\usepackage[T1]{fontenc}
\usepackage[utf8]{inputenc}
\usepackage{lmodern}
\usepackage{geometry}
\geometry{margin=27mm}
\usepackage{amsmath,amssymb,amsthm,mathtools}
\usepackage{booktabs,longtable,array}
\usepackage{enumitem}
\usepackage{xcolor}
\usepackage{hyperref}
\hypersetup{colorlinks=true,linkcolor=blue!45!black,urlcolor=blue!45!black,citecolor=blue!45!black}
\usepackage{microtype}

\newtheorem{theorem}{Theorem}[section]
\newtheorem{proposition}[theorem]{Proposition}
\newtheorem{lemma}[theorem]{Lemma}
\newtheorem{corollary}[theorem]{Corollary}
\theoremstyle{definition}
\newtheorem{definition}[theorem]{Definition}
\newtheorem{conjecture}[theorem]{Conjecture}
\newtheorem{example}[theorem]{Example}
\theoremstyle{remark}
\newtheorem{remark}[theorem]{Remark}
\newtheorem{warning}[theorem]{Warning}

\newcommand{\SN}{S_N}
\newcommand{\TN}{T_N}
\newcommand{\sm}{\operatorname{sum}}
\newcommand{\RR}{\mathcal{R}}
\DeclareMathOperator{\val}{val}

\newcommand{\notesep}{\par\smallskip\noindent}

\title{\textbf{Catch-Up on Consecutive Integers}\\[3pt]
\large Precise formulation, rigorous partial results,\\
proof attempts, and a counterexample to a proposed strategy}
\author{GPT 6 Astra Pro}
\date{8 October 2026}

\begin{document}
\noindent\textbf{Problem ID:} \texttt{C72-14}\quad\textbf{Model:} GPT 6 Astra Pro\quad\textbf{Version:} 1\par
\noindent\textbf{Status:} Unresolved for the complete catalogue problem; this manuscript records partial results and research attempts.\par
\maketitle

\begin{abstract}
For Catch-Up played on $\{1,2,\dots,N\}$, the conjecture is that
optimal play produces a draw whenever $N(N+1)/2$ is even.  These notes
collect the proof attempts made toward that statement.  We give a
precise sequential model, an exact minimax recurrence, an opening
antisymmetry (an observation attributed to Bernhard von Stengel), a
sufficient two-opening criterion, and a proof that every even-total
consecutive instance admits \emph{some} legal drawn play.  We then
formulate half-score completion and exchange conditions, prove a
one-deletion repair property of the canonical four-number blocks,
and isolate the defect in a proposed small-deficit invariant.  An
explicit adversarial analysis for $N=8$ proves that a strategy forced
to maintain this small-deficit invariant loses, although an excluded
move is drawing.  We distinguish valid lemmas, equivalent
reformulations, sufficient conditions, and failed strategies.
No uniform winning-or-drawing strategy for all $N$ is obtained.
\end{abstract}

\tableofcontents


\section*{Formal problem statement: Catch-Up on consecutive integers}
\addcontentsline{toc}{section}{Formal problem statement: Catch-Up on consecutive integers}
\noindent
\textbf{Instance.} Fix an integer $N\ge 1$ and the finite set
$S_N=\{1,2,\ldots,N\}$. Two players, $P_1$ and $P_2$, start with
scores $s_1=s_2=0$. Player $P_1$ makes the first move by taking exactly
one integer. Thereafter the designated player takes unused integers one
at a time, and is required to keep taking integers while their score is
strictly below the other player's score; the turn ends immediately upon
equality or overtaking. If the player cannot catch up even by taking all
remaining integers, that player takes every remaining integer and the
game terminates. Play also terminates when all integers have been taken.
The terminal payoff to $P_1$ is the final score difference $s_1-s_2$.
Both players have perfect information and choose optimally, with $P_1$
maximising this payoff and $P_2$ minimising it.

\medskip\noindent
\textbf{Conjectured answer.} Put $T_N=N(N+1)/2$. For every $N$ for
which $T_N$ is even (equivalently, $N\equiv0,3\pmod4$), prove or
disprove that the optimal terminal score difference is exactly zero:
\[
\boxed{\operatorname{Val}(\operatorname{Catch\mbox{-}Up}(S_N))=0.}
\]
Equivalently, under optimal play both players finish with score
$T_N/2=N(N+1)/4$. The existence of a legal drawn sequence of moves
is weaker than the claim that the minimax outcome is a draw.


\section{Rules and the question}

\begin{definition}[Catch-Up]
Let $S$ be a finite set of \emph{distinct positive integers}.  Players
$P_1$ and $P_2$ begin with respective scores $(s_1,s_2)=(0,0)$.
The following rules apply.
\begin{enumerate}[label=(\roman*)]
\item $P_1$ starts by choosing \emph{exactly one} element of $S$,
      removing it, and adding it to her score.
\item Thereafter, the player designated to act selects remaining
      elements \emph{one at a time}, adding each to that player's score.
      The acting player must continue selecting while strictly behind,
      and must stop as soon as their score \emph{equals or exceeds} the
      opponent's current score.  The turn then passes to the opponent,
      who must make at least one selection if the remaining set is
      nonempty, including if the scores are tied.
\item If the acting player is behind and cannot catch up even by taking
      every remaining number, that player receives all remaining
      numbers and the game ends.  This convention agrees with
      continuing to select the numbers one by one until none remain.
\item When no numbers remain, the player with the larger score wins;
      equality means a draw.
\end{enumerate}
Every selected number is allocated to exactly one player; thus the
sum of the final scores equals $\sm(S):=\sum_{x\in S}x$.
\end{definition}

\begin{example}[A drawn play for $N=4$]
From $\{1,2,3,4\}$, let $P_1$ choose $2$.  Player $P_2$ chooses $1$
(still behind) and then $4$ (now ahead), and $P_1$ takes $3$.
The final scores are $(5,5)$.  The existence of this play does not
by itself prove a draw under optimal play.
\end{example}

Write
\[
\SN=\{1,2,\ldots,N\},\qquad
\TN=\sm(\SN)=\frac{N(N+1)}2.
\]
When $\TN$ is even, put $H=\TN/2$.  The problem is the following.

\begin{conjecture}[Optimal-draw conjecture for consecutive integers]
\label{conj:main}
For every positive integer $N$ with
\[
N\equiv 0\text{ or }3\pmod 4,
\]
optimal play in $\operatorname{Catch\mbox{-}Up}(\SN)$ ends in a draw;
that is, both final scores are $H=N(N+1)/4$.
\end{conjecture}

\begin{remark}[Meaning of optimal play]
The game is finite, has perfect information and no chance moves.
We use the final difference $s_1-s_2$ as a zero-sum payoff: $P_1$
maximises it and $P_2$ minimises it.  This is a convenient refinement
of preferences among wins or losses.  In particular, the sign of the
minimax payoff gives the same win--draw--loss outcome as the ordinary
preferences win $>$ draw $>$ loss.  The conjecture asserts that the
minimax payoff is zero; it is not an assertion that every legal play
is drawn.
\end{remark}

\section{Exact game-value recurrence}

\begin{definition}[The continuation value]
Let $V(R,d)$ be the minimax \emph{final score advantage of the player
about to select a number}, when the remaining set is $R$ and that
player is currently behind by $d\ge0$.  When $d=0$ and $R\ne\varnothing$,
the designated player is nevertheless required to select a number.
Write $\sm(R)=\sum_{x\in R}x$, with $\sm(\varnothing)=0$.
\end{definition}

\begin{proposition}[Exact one-selection recurrence]
\label{prop:recurrence}
For every $d\ge0$,
\begin{equation}\label{eq:terminal}
V(\varnothing,d)=-d.
\end{equation}
For nonempty $R$,
\begin{equation}\label{eq:recursion}
\boxed{
V(R,d)=\max_{x\in R}
\begin{cases}
V(R\setminus\{x\},d-x),&x<d,\\[2pt]
-\,V(R\setminus\{x\},x-d),&x\ge d.
\end{cases}}
\end{equation}
If $\sm(R)<d$, then $V(R,d)=\sm(R)-d$.
\end{proposition}
\begin{proof}
If $x<d$, the acting player is still behind after taking $x$,
so that same player must select again; the residual deficit is $d-x$.
If $x\ge d$, the player has caught up or overtaken, and the next
selection belongs to the opponent, now behind by $x-d$.  This
interchanges the player whose advantage is measured, producing the
minus sign.  In particular, \emph{equality $x=d$ belongs in the
second branch}: a tied position passes the turn.  Maximisation over
$x$ expresses the acting player's choice.  If $R=\varnothing$, no
more points can be scored and the designated player's margin is $-d$.
Finally, if $\sm(R)<d$, the acting player cannot catch up and receives
every remaining number, giving margin $\sm(R)-d$.
\end{proof}

Applying the recurrence to the initial tied position gives
\begin{equation}\label{eq:openingvalue}
V(\SN,0)=-\min_{x\in\SN}V(\SN\setminus\{x\},x).
\end{equation}
Consequently, Conjecture~\ref{conj:main} is equivalent to
\begin{align}
V(\SN\setminus\{x\},x)&\ge0
          &&\text{for \emph{every} }x\in\SN,\label{eq:all-opening}\\
V(\SN\setminus\{x\},x)&\le0
          &&\text{for \emph{at least one} }x\in\SN.\label{eq:one-opening}
\end{align}
The first line is the assertion that $P_2$ can avoid losing against
any first move; the second that $P_1$ has at least one non-losing
opening.  Both directions are needed.

\section{Opening reversal and the two-opening criterion}

\noindent
The opening-reversal observation and its use are attributed here to
Bernhard von Stengel, as indicated by the problem proposer; they are
not asserted to be new results of these notes.

For distinct $a,b\in S$, define $M_S(a,b)$ as the optimal final
margin $s_1-s_2$ conditional on $P_1$ first taking $a$ and $P_2$
first taking $b$, with optimal continuation thereafter.  This is a
conditional value after the \emph{first selections} of the players;
if $b<a$, $P_2$ has not yet completed his turn.

\begin{lemma}[Opening reversal]
\label{lem:reverse}
For any finite $S$ of distinct positive numbers and $a\ne b$ in $S$,
\begin{equation}\label{eq:antisym}
\boxed{M_S(a,b)=-M_S(b,a).}
\end{equation}
More precisely, for $a<b$,
\[
M_S(a,b)=V(S\setminus\{a,b\},b-a),
\qquad
M_S(b,a)=-V(S\setminus\{a,b\},b-a).
\]
\end{lemma}
\begin{proof}
Assume $a<b$.  After $P_1$ takes $a$ and $P_2$ takes $b$,
$P_2$ has overtaken and $P_1$ is about to select from
$S\setminus\{a,b\}$, behind by $b-a$.  Its value, measured for $P_1$,
is $V(S\setminus\{a,b\},b-a)$.  In the reversed opening, $P_1$
takes $b$ and $P_2$ takes $a$.  Player $P_2$ is still behind and must
continue the \emph{same} turn; the remaining set and deficit are
identical, but the acting player's identity is reversed.  Therefore
the values have opposite signs.
\end{proof}

\begin{proposition}[Opening form of the minimax value]
\label{prop:matrix}
For a set $S$ with at least two elements,
\begin{equation}\label{eq:matrix}
\boxed{V(S,0)=\max_{a\in S}\ \min_{b\in S\setminus\{a\}}M_S(a,b).}
\end{equation}
\end{proposition}
\begin{proof}
Player $P_1$ chooses an opening $a$, and $P_2$ then chooses a first
number $b\ne a$.  Further selections by $P_2$, if necessary, and
all later turns are optimised in $M_S(a,b)$.  Thus $P_2$ minimises
over $b$ and $P_1$ maximises over $a$.
\end{proof}

An opening $a$ is called \emph{non-losing for $P_1$} when
$\min_{b\ne a}M_S(a,b)\ge0$.

\begin{corollary}[At most one strictly winning opening]
\label{cor:unique}
There cannot be two distinct openings from which $P_1$ can force a
strictly positive final margin.
\end{corollary}
\begin{proof}
If $a$ is a strictly winning opening, then $M_S(a,b)>0$ for
all $b\ne a$.  By Lemma~\ref{lem:reverse}, $M_S(b,a)<0$ for every
$b\ne a$, and therefore every other opening has a losing reply.
\end{proof}

\begin{theorem}[Two non-losing openings suffice for a draw]
\label{thm:twoopen}
If there exist two distinct non-losing openings $a,b\in S$, then
$V(S,0)=0$.
\end{theorem}
\begin{proof}
The existence of $a$ already gives $V(S,0)\ge0$.  Moreover,
$M_S(a,b)\ge0$ and $M_S(b,a)\ge0$, so antisymmetry forces
$M_S(a,b)=M_S(b,a)=0$.  Player $P_2$ can avoid losing as follows:
against any opening $x\ne a$, select $a$ first, obtaining
\[
M_S(x,a)=-M_S(a,x)\le0.
\]
Against opening $a$, select $b$ first, yielding $M_S(a,b)=0$.
Hence $P_2$ can hold the value to at most zero.  Together with
$V(S,0)\ge0$, this gives $V(S,0)=0$.
\end{proof}

\begin{remark}[The remaining gap]
To turn Theorem~\ref{thm:twoopen} into a proof of
Conjecture~\ref{conj:main}, one still needs a \emph{uniform
construction of two non-losing openings} for every $N\equiv0,3
\pmod4$.  Neither the opening symmetry nor equal-sum arithmetic
supplies that construction.  Antisymmetry by itself does not imply
that any row of the matrix $M_S$ is nonnegative.
\end{remark}

\section{Equal-sum partitions: what they actually prove}

\begin{proposition}[Parity]
\label{prop:parity}
For $N\ge1$,
\[
\TN\text{ is even}\quad\Longleftrightarrow\quad
N\equiv0\text{ or }3\pmod4.
\]
\end{proposition}
\begin{proof}
The product $N(N+1)$ is divisible by $4$ exactly when one of the
consecutive integers $N,N+1$ is divisible by $4$.
\end{proof}

\begin{lemma}[Explicit equal-sum partition]
\label{lem:partition}
Whenever $\TN$ is even, $\SN$ admits a disjoint partition
$\SN=A\sqcup B$ with
\[
\sm(A)=\sm(B)=H=\frac{\TN}{2}.
\]
\end{lemma}
\begin{proof}
If $N=4m$, split each block
$\{4j-3,4j-2,4j-1,4j\}$, $1\le j\le m$, as
\[
\{4j-3,4j\}\quad\text{and}\quad\{4j-2,4j-1\}.
\]
Both pairs have sum $8j-3$.  Assign one pair from each block to
$A$, the other to $B$.

If $N=4m+3$, first split $\{1,2,3\}$ into $\{3\}$ and
$\{1,2\}$, each summing to $3$.  Then for $1\le j\le m$ split
$\{4j,4j+1,4j+2,4j+3\}$ into
\[
\{4j,4j+3\}\quad\text{and}\quad\{4j+1,4j+2\},
\]
whose sums are both $8j+3$.  Combining the assigned halves produces
the claimed partition.
\end{proof}

\begin{theorem}[A legal drawn play exists]
\label{thm:exists}
Let $S$ be any finite set of distinct positive integers that has an
equal-sum partition $S=A\sqcup B$, with
$\sm(A)=\sm(B)=H$.  Then \emph{there exists a legal play} of
Catch-Up$(S)$ that ends in a draw, with final scores $(H,H)$.
In particular, this is true for $S=\SN$ whenever $\TN$ is even.
\end{theorem}
\begin{proof}
Give $A$ to $P_1$ and $B$ to $P_2$ as \emph{voluntary selection
sets}.  Each player selects only numbers from their own assigned
set, one by one, stopping at the first selection on a turn that
catches or overtakes the opponent.

Suppose the acting player has score $u$ and the opponent has score
$v$, where $u\le v\le H$.  The acting player still has assigned
numbers with total sum $H-u$.  Since
\[
H-u\ge v-u,
\]
there is enough assigned weight to catch up; selecting assigned
numbers in any order and stopping at the first catch-up is legal.
When the scores are tied below $H$, the same player has an assigned
number left and can start the next turn.  Every score remains at
most $H$, because nobody selects outside their own assigned set.
If a player reaches $H$, the other player receives any remaining
numbers and reaches $H$ as well.  Otherwise play continues until both
assigned sets are exhausted, again giving $(H,H)$.
\end{proof}

\begin{warning}[Arithmetic feasibility is not strategic force]
Theorem~\ref{thm:exists} assumes both players \emph{cooperate} with
the assignment.  An opponent is free to take numbers from the other
player's intended half.  Therefore this theorem does \emph{not}
imply that either player can force a draw under optimal play.
\end{warning}

\begin{example}[Naive complementary pairing fails]
For $S=\{1,2,3,4\}$, the partition
$\{1,4\}\sqcup\{2,3\}$ has equal sums.  A proposed response
``answer an opening $1$ with its partner $4$'' loses: the legal
sequence $P_1:1$, $P_2:4$, $P_1:2,3$ ends at $(6,4)$.
Thus a fixed-pair reply is not itself a non-losing strategy.
\end{example}

\section{Half-score completion sets and the failed induction shortcut}

\begin{lemma}[Reaching the half-total]
\label{lem:half}
Suppose the sum of all initial numbers is $2H$.
If either player's score first reaches \emph{exactly} $H$, the final
outcome is forced to be a draw.  If either player's score ever exceeds
$H$, that player has already secured a win.
\end{lemma}
\begin{proof}
When a player first reaches $H$, that player's turn has ended,
because the opponent cannot have more than $H$ (the total is $2H$).
If the opponent's current score is $t\le H$, the remaining numbers
have sum $H-t$.  The opponent, now acting, must receive all of them
in order to catch up, so the final scores are $(H,H)$.
If a player instead exceeds $H$, the sum of every number allocated
to the opponent is strictly less than $H$.  The player who exceeded
$H$ therefore necessarily wins.
\end{proof}

\begin{definition}[Half-score completion set]
At a position with total initial sum $2H$, a \emph{completion set}
for a player with score $s\le H$ is a subset $A$ of the currently
remaining numbers such that $\sm(A)=H-s$.
\end{definition}

\begin{remark}[Why completion sets alone do not solve the game]
Having a completion set gives a player enough designated numbers to
reach $H$ while the opponent's score is at most $H$.  But the
opponent may take elements of that set on intervening turns.
Preservation against an adversary, not initial existence, is the
missing ingredient.
\end{remark}

\begin{example}[A residual tied game need not retain its value]
\label{ex:residual}
The set $R=\{3,4,7\}$ has total $14$, and the fresh game on $R$
is drawn under optimal play: the first player can take $7$, forcing
the second player to take $3$ and $4$; conversely, a second player
can reply to an opening $3$ or $4$ by taking $7$, resulting in a draw.
Thus
\[
V(\{3,4,7\},0)=0.
\]
But for an acting player \emph{already behind by $2$},
\begin{equation}\label{eq:offsetexample}
V(\{3,4,7\},2)=-2.
\end{equation}
Indeed, taking $3$ leaves $\{4,7\}$; the opponent chooses $7$ and
the acting player eventually receives $4$.  Taking $4$ is symmetric:
the opponent takes $7$ and the player eventually receives $3$.
Taking $7$ makes the opponent take both $3$ and $4$ to catch up.
In every case the original acting player loses by $2$.

This offset position is reachable within a \emph{legal} play on
$\{1,\ldots,7\}$:
\[
P_1:1,\qquad P_2:2,\qquad P_1:5,\qquad P_2:6,
\]
which leaves scores $(6,8)$ and $R=\{3,4,7\}$ with $P_1$ to act.
This is not a counterexample to Conjecture~\ref{conj:main}; the
prefix is not claimed to be optimal.
\end{example}

\begin{warning}[Invalid induction shortcut]
An induction on $N$ cannot replace the remaining game by a \emph{fresh}
Catch-Up game merely because its remaining numbers admit an equal-sum
partition.  The score deficit $d$ is part of the position, and
Example~\ref{ex:residual} shows that changing $d$ changes the
optimal value.  Any induction must track both the remaining set
and the current deficit.
\end{warning}

\section{Balanced completions and signed exchanges}

Assume a position has current scores $a$ and $a+d$, with $d\ge0$,
and the player at $a$ about to select.  Set
\[
r=\sm(R),\qquad 2a+d+r=2H,\qquad
q=H-a=\frac{r+d}{2}.
\]
The acting player needs $q$ points to reach $H$; the opponent
needs $q-d$.

\begin{definition}[Balanced completion partition]
A partition $R=A\sqcup B$ is \emph{balanced} for this position when
\begin{equation}\label{eq:balanced}
\boxed{\sm(A)=q,\qquad\sm(B)=q-d.}
\end{equation}
The sets $A$ and $B$ complete the acting and nonacting players,
respectively, to $H$.
\end{definition}

\begin{lemma}[Completion-partition equivalence]
\label{lem:equiv}
A balanced completion partition exists if and only if the acting
player has a half-score completion subset $A\subseteq R$.
\end{lemma}
\begin{proof}
One implication is immediate.  Conversely, if $\sm(A)=H-a=q$,
then its complement has sum
\[
\sm(R\setminus A)=r-q
=(2H-2a-d)-(H-a)=H-a-d=q-d.
\]
Thus $B=R\setminus A$ gives the balanced partition.
\end{proof}

\begin{remark}[Correction to an earlier proposed strengthening]
Using two labelled halves is useful for describing exchanges, but
Lemma~\ref{lem:equiv} shows that \emph{balanced completion} is not
logically stronger than the original single completion-set condition.
\end{remark}

\begin{lemma}[Propagation under cooperation]
\label{lem:propagate}
Assume $R=A\sqcup B$ satisfies \eqref{eq:balanced}.  If the acting
player legally completes a turn using only a subset $C\subseteq A$
with $\sm(C)=d+e$, where $e\ge0$ is the new lead, then the new
remaining set $(A\setminus C)\sqcup B$ is balanced for the other
player's next turn, with the halves exchanged.
\end{lemma}
\begin{proof}
After the turn, the other player is behind by $e$ and still needs
$q-d$ to reach $H$.  Its designated completion set is $B$, of sum
$q-d$.  The former acting player's remaining half has sum
\[
\sm(A\setminus C)=q-(d+e)=q-d-e,
\]
which is precisely the former acting player's remaining distance
from $H$.  The asserted balanced partition follows.
\end{proof}

The essential issue is that an opponent may select numbers from
\emph{both} halves.  The next identity records the exact repair
requirement without asserting that it can always be met.

\begin{lemma}[Exact signed-exchange equation]
\label{lem:exchange}
Suppose $R=A\sqcup B$ satisfies \eqref{eq:balanced}, and the acting
player makes selections comprising a set $C\subseteq R$.  Write
\[
C_A=C\cap A,\qquad C_B=C\cap B,\qquad w=\sm(C_B).
\]
The other player, whose score was $a+d$ and is unchanged during this
turn, still needs $q-d$ points to reach $H$.  Choose subsets
\[
D\subseteq A\setminus C_A,\qquad
E\subseteq B\setminus C_B
\]
and let
\[
B'=\bigl((B\setminus C_B)\setminus E\bigr)\cup D.
\]
Then $B'$ is a completion set for the other player \emph{if and only if}
\begin{equation}\label{eq:exchange}
\boxed{\sm(D)-\sm(E)=w.}
\end{equation}
\end{lemma}
\begin{proof}
Since $\sm(B)=q-d$,
\[
\sm(B')=(q-d)-w-\sm(E)+\sm(D).
\]
This equals the required $q-d$ precisely when
$\sm(D)-\sm(E)=w$.
\end{proof}

\begin{remark}[A necessary equation, not an existence theorem]
The signed exchange \eqref{eq:exchange} is less restrictive than
requiring a subset of $A\setminus C_A$ to sum exactly to $w$:
we may also give up some elements of $B\setminus C_B$.
Nevertheless, writing down the equation does not prove that a
solution $(D,E)$ exists.  In particular, a hypothetical move that
makes the acting player's score exceed $H$ leaves the opponent
unable to complete to $H$ at all.
\end{remark}

\section{Four-number block repairs and a proposed invariant}

\begin{lemma}[Repair after one deletion in an intact block]
\label{lem:block}
Let $u$ be a positive integer and
\[
Q=\{u,u+1,u+2,u+3\},\qquad
O=\{u,u+3\},\qquad I=\{u+1,u+2\}.
\]
Then $\sm(O)=\sm(I)=2u+3$.  After deleting \emph{any one}
number $x\in Q$, the remaining three numbers contain a subset of
sum $2u+3$.
\end{lemma}
\begin{proof}
If $x\in O$, the set $I$ survives and has sum $2u+3$.
If $x\in I$, the set $O$ survives and has the same sum.
\end{proof}

This lemma motivates reorienting a block after the opponent
``steals'' a number.  For example, if $O$ was the selected completion
half and $u$ was removed, replacing the surviving $\{u+3\}$ by
$I$ increases the completion-side sum by exactly $u$:
\[
\sm(I)-(u+3)=(2u+3)-(u+3)=u.
\]
This is genuine local flexibility, but applies to an \emph{intact}
block before a \emph{single} deletion.  It says nothing by itself
about repeated hits to a partially consumed block, let alone the
legality of a complete catch-up move.

\begin{lemma}[Small deficit: one selection and no crossing of $H$]
\label{lem:small}
Suppose a position has a balanced completion partition
$R=A\sqcup B$, with the acting player behind by $d\ge0$, and assume
\begin{equation}\label{eq:small}
\boxed{d\le\min R.}
\end{equation}
Then the acting player must complete the turn in \emph{exactly one}
selection and cannot finish that turn with score greater than $H$.
If the player reaches $H$, the game is drawn.
\end{lemma}
\begin{proof}
Every $x\in R$ satisfies $x\ge d$, so the first selection catches
up or overtakes and the turn immediately ends.  By
\eqref{eq:balanced}, any $x\in R$ lies in a subset whose sum is at
most $q$: indeed $x\le\sm(A)=q$ if $x\in A$, and
$x\le\sm(B)=q-d\le q$ if $x\in B$.  The acting player's score
becomes $a+x\le a+q=H$.  Equality invokes Lemma~\ref{lem:half}.
\end{proof}

The attempted induction aimed to maintain \emph{both} a balanced
completion partition and \eqref{eq:small} immediately after every
turn of one favoured player.  The opposing player would then get
only one choice before the next repair.  The local repair lemma makes
this tempting.  To complete such a programme one would also need
the following \emph{sufficient move criterion}.

\begin{lemma}[A sufficient small-overshoot move criterion]
\label{lem:overshoot}
Suppose a player is about to select from $R$ and is behind by
$e>0$.  Suppose $C\subseteq R$ is nonempty, has sum
$\sm(C)=e+e'$ for some $e'\ge0$, and contains a number $y>e'$.
Then the player may select all elements of $C\setminus\{y\}$ in
any order and $y$ last, thereby ending the turn \emph{legally}
with lead exactly $e'$.  If moreover $R\setminus C\ne\varnothing$ and
\[
e'\le\min(R\setminus C),
\]
then the resulting handoff satisfies the small-deficit bound.

When the players are initially tied ($e=0$), a legal turn consists
of \emph{exactly one} selection $y$; it leaves lead $e'=y$ and
satisfies the small-deficit bound only if
$R\setminus\{y\}\ne\varnothing$ and
$y\le\min(R\setminus\{y\})$.
\end{lemma}
\begin{proof}
For $e>0$, just before the last selection the acting player has
added
\[
\sm(C)-y=e+e'-y<e.
\]
Because every number is positive, each earlier partial sum is also
strictly below $e$.  Hence no earlier selection would have ended
the turn.  The final choice $y$ takes the cumulative addition to
$e+e'\ge e$, which ends the turn with lead $e'$.  The additional
inequality is exactly the desired small-deficit condition for the
next player.  At a tie, every positive first selection immediately
creates a lead and ends the turn, so a multiple-selection turn is
impossible; the separate statement follows.
\end{proof}

\begin{remark}[Corrections and limitations]
An earlier formulation of the sufficient criterion implicitly
required a final element $y>e'$ even when $e=0$.  That would be
impossible because $\sm(C)=e'$ at a tie.  The tied case must be
treated separately as in Lemma~\ref{lem:overshoot}.  More importantly,
this lemma does \emph{not} assert that the needed subset $C$ exists.
Even if suitable blocks admit exchanges, imposing the small-deficit
condition turns out to exclude essential drawing moves.
\end{remark}

\section{Rigorous obstruction at $N=8$}

In this section, the \emph{small-deficit restriction} means that a
specified player must, after each of their completed turns and
before a terminal outcome, leave a remaining set $R$ and score lead
$\ell\ge0$ satisfying
\begin{equation}\label{eq:A}
\tag{A}\boxed{\ell\le\min R.}
\end{equation}
No completion-partition condition is imposed: the restriction is
tested \emph{in isolation}, making the negative result stronger.

\begin{proposition}[Small-deficit strategies lose for $N=8$]
\label{prop:N8}
In Catch-Up$(\{1,\ldots,8\})$, $P_2$ has a strategy defeating
\emph{every} $P_1$ strategy constrained to satisfy (A) after each of
her turns until the game finishes.  In particular, (A) is not a
valid invariant for a uniform drawing strategy.
\end{proposition}
\begin{proof}
The total is $36$, so $H=18$.  At her first selection, (A) forces
$P_1$ to take $1$: if she took any $x>1$, her lead $x$ would exceed
the still-available minimum $1$.  Player $P_2$ responds with $2$,
leaving
\[
(s_1,s_2)=(1,2),\qquad R=\{3,4,5,6,7,8\}.
\]
Player $P_1$ is behind by $1$ and therefore obtains exactly one
selection $x$.  Its resulting lead is $x-1$.  Condition (A) allows
only $x=3$ or $x=4$: all $x\ge5$ leave $3$ available and give
$x-1>3$.

\smallskip\noindent\textbf{Case I: $P_1$ takes $3$.}
Player $P_2$ answers with $4$, giving
\[
(s_1,s_2)=(4,6),\qquad R=\{5,6,7,8\}.
\]
Now $P_1$ is behind by $2$, so she must again take one number.
Taking $8$ would create lead $4+8-6=6$ while leaving $5$
available, violating (A).  Thus $P_1$ must take
$x\in\{5,6,7\}$.  Player $P_2$ responds with $8$, leaving
\[
(s_1,s_2)=(4+x,14),\qquad
R=\{5,6,7\}\setminus\{x\}.
\]
Each remaining number is at least the new deficit $10-x$ of $P_1$.
So she takes one of the two remaining numbers, say $z$, and $P_2$
takes the final number $w$.  Consequently,
\[
s_1=4+x+z\le17,
\qquad
s_2=14+w\ge19.
\]
Thus $P_1$ loses in every permitted continuation.

\smallskip\noindent\textbf{Case II: $P_1$ takes $4$.}
Player $P_2$ responds with $5$, obtaining
\[
(s_1,s_2)=(5,7),\qquad R=\{3,6,7,8\}.
\]
Player $P_1$, behind by $2$, must take one number.  Taking
$x\in\{6,7,8\}$ would give lead $x-2>3$ while leaving $3$,
violating (A).  Hence she must choose $3$, reaching $(8,7)$.
Player $P_2$ takes $7$, reaching
\[
(s_1,s_2)=(8,14),\qquad R=\{6,8\}.
\]
There are two legal endings:
\begin{center}
\begin{tabular}{@{}ccc@{}}
\toprule
$P_1$ takes & $P_2$ takes & Final $(s_1,s_2)$\\
\midrule
$6$ & $8$ & $(14,22)$\\
$8$ & $6$ & $(16,20)$\\
\bottomrule
\end{tabular}
\end{center}
Both are losses for $P_1$.  This completes the proof.
\end{proof}

The next proposition confirms that the restriction excludes a
\emph{genuinely drawing} continuation rather than merely some
unneeded, aesthetically unpleasant choice.

\begin{proposition}[An excluded move is the only non-losing move]
\label{prop:drawing8}
At the legal position
\begin{equation}\label{eq:critical}
\boxed{(s_1,s_2)=(4,6),\qquad R=\{5,6,7,8\},
\qquad P_1\text{ to move},}
\end{equation}
$P_1$ can avoid losing by taking $8$, whereas each of the choices
$5,6,7$ is losing against the response $8$ by $P_2$.
The drawing move $8$ violates (A), although it leaves a balanced
completion partition.
\end{proposition}
\begin{proof}
The losses after selecting $5,6,$ or $7$ are established in
Case I of Proposition~\ref{prop:N8}.

If instead $P_1$ takes $8$, the scores become $(12,6)$ and
$R=\{5,6,7\}$.  Player $P_2$ is behind by $6$.
His \emph{complete legal turns} are precisely
\[
\langle6\rangle,\quad\langle7\rangle,\quad
\langle5,6\rangle,\quad\langle5,7\rangle.
\]
For each, $P_1$ can respond as follows:
\begin{center}
\begin{tabular}{@{}ccc@{}}
\toprule
$P_2$'s turn & $P_1$ next takes & Final $(s_1,s_2)$\\
\midrule
$\langle6\rangle$ & $7$ & $(19,17)$\\
$\langle7\rangle$ & $6$ & $(18,18)$\\
$\langle5,6\rangle$ & $7$ & $(19,17)$\\
$\langle5,7\rangle$ & $6$ & $(18,18)$\\
\bottomrule
\end{tabular}
\end{center}
Thus $P_1$ can force at least a draw; the second or fourth reply
shows that $P_2$ can force a draw against this choice, so the move's
minimax outcome is exactly a draw.

Finally, after taking $8$, $P_1$ leads by $6$, whereas the
smallest remaining number is $5$.  Hence (A) fails.  Nevertheless,
$P_1$ needs $18-12=6$ to reach $H$, and $P_2$ needs $18-6=12$;
the remaining set has balanced completion partition
\[
\underbrace{\{6\}}_{\sm=6}
\qquad\sqcup\qquad
\underbrace{\{5,7\}}_{\sm=12}.
\]
So balanced completion survives even though small deficit does not.
\end{proof}

\begin{corollary}[Failure of the proposed proof scheme]
\label{cor:failure}
There is no proof of Conjecture~\ref{conj:main} based on the claim
that $P_1$ can always play non-losing while maintaining (A) after
every turn: that claim is false already for $N=8$.
\end{corollary}
\begin{proof}
Proposition~\ref{prop:N8} directly contradicts the claimed
strategy.  Adding a completion-partition requirement or further
blockwise rules cannot restore a strategy while retaining the
restriction (A).
\end{proof}

\section{What has and has not been established}

\subsection*{Proved results}
\begin{enumerate}[label=\arabic*.]
\item The minimax recurrence \eqref{eq:recursion} exactly encodes the
one-by-one selection rule, including turn changes upon equality.
\item Opening reversal, the criterion of two non-losing openings,
and uniqueness of a strictly winning opening are valid; the opening
symmetry is attributed to Bernhard von Stengel.
\item For $N\equiv0,3\pmod4$, an explicit equal-sum partition exists
and yields at least one legal drawn play.
\item Reaching the half-total exactly forces a draw; exceeding it
forces a win.
\item The single completion-subset condition is equivalent to the
balanced completion-partition condition; repair by a signed exchange
is governed exactly by \eqref{eq:exchange}.
\item An intact four-number block can retain its half-block target
sum after one arbitrary deletion, but this does not establish a
repeated repair strategy.
\item The combined small-deficit strategy fails, even without its
partition component, by an explicit $N=8$ adversarial analysis.
\end{enumerate}

\subsection*{Claims \emph{not} established}
\begin{enumerate}[label=\arabic*.]
\item A uniform construction of two non-losing openings for
$\SN$ whenever $\TN$ is even.
\item Preservation of half-score completion sets under all legal
adversarial turns.
\item A repair-and-catch-up exchange strategy across repeatedly used
four-number blocks.
\item An induction step from $N$ to $N+4$ preserving the full game
value when the residual deficit is included.
\item Conjecture~\ref{conj:main}, i.e. the uniform assertion
$V(\SN,0)=0$ for every $N\equiv0,3\pmod4$.
\end{enumerate}

\begin{remark}[A more appropriate future direction]
The example in Proposition~\ref{prop:drawing8} shows that a drawing
strategy must sometimes \emph{permit} an opponent to make multiple
selections on the next turn.  A plausible invariant would need to
track the full residual game state, or classify safe multi-selection
responses while permitting completion-set repairs.  This is only a
research direction; no such invariant or induction is proved here.
\end{remark}

\section*{Notes on attribution and computational evidence}
\addcontentsline{toc}{section}{Notes on attribution and computational evidence}
The original Catch-Up article is due to Isaksen, Ismail, Brams and
Nealen (2015)~\cite{isaksen2015}.  The author of the present problem
attributes the opening-reversal result to Bernhard von Stengel; its
inclusion above is expository rather than a priority claim.
Von Stengel's conference abstract reports computational results
through $N=28$ and discusses the desirability of an induction
proof~\cite{vstengel2019}.  The problem poser reports additional
brute-force checks of the consecutive even-total case through
$N=34$; those checks are \emph{not} reproduced or independently
verified in these notes.  Finite computational verification,
regardless of its range, does not substitute for a general proof.

\begin{thebibliography}{9}
\bibitem{isaksen2015}
A.~Isaksen, M.~Ismail, S.~J.~Brams, and A.~Nealen,
\emph{Catch-Up: A Game in Which the Lead Alternates},
Game \& Puzzle Design, vol.~1, no.~2 (2015), pp.~38--49.
\url{https://mpra.ub.uni-muenchen.de/108784/}.

\bibitem{vstengel2019}
B.~von Stengel, \emph{Computational progress on the Catch-Up game},
conference abstract in CGTC III, Associa\c{c}\~ao Ludus.
\url{https://ludicum.org/en/cgtc-iii/}.

\bibitem{formalconjectures}
\emph{The Catch-Up Game Conjecture}, formal-conjectures project,
issue \#1324.
\url{https://github.com/google-deepmind/formal-conjectures/issues/1324}.
\end{thebibliography}
\section{Completion Estimate}
10\%

\end{document}
